%% kapa: exam=1 round=37 year=2026 period=1 form=A subject=real_estate subject_ko=부동산학원론 seq=26 number=106
%% source: Q-Net 감정평가사 기출문제 (cst00309) artlSeq=5258924 file=raw/1008/5258924/2026년 제37회 감정평가사 1차 1교시.pdf
%% conversion: kordoc 4.24.51 -> tools/gen1.py (auto); review=proofread-sonnet; flags=none
\begin{question}{106}
부동산증권에 관한 설명으로 옳은 것은?
\begin{choices}
  \item 저당이체증권(MPTS)의 투자자는 금리위험과 채무불이행위험을 부담한다.
  \item 저당이체채권(MPTB)의 발행자는 조기상환위험을 부담하며 투자자는 채무불이행위험을 부담한다.
  \item 저당담보부채권(MBB)의 투자자는 금리위험과 채무불이행위험을 부담한다.
  \item 다른 조건이 일정할 때 시장금리가 표면금리보다 높은 경우 저당담보부채권(MBB)의 발행가격은 높아진다.
  \item 다른 조건이 일정할 때 저당이체채권(MPTB)의 듀레이션(duration)은 조기상환이 많이 일어난 경우가 조기상환이 없는 경우보다 길다.
\end{choices}
\end{question}
\answer{1}
